% ============================================================
% SECCIÓN 1: INTRODUCCIÓN
% ============================================================

\section{Introducción}

El 7 de mayo de 2016, un Tesla Model S con el sistema Autopilot activo colisionó fatalmente con un tractor-tráiler en Williston, Florida. La investigación oficial de la NHTSA concluyó que no hubo defecto: el sistema funcionó exactamente como fue diseñado \citep{nhtsa2017pe16007}. El sistema de frenado automático nunca fue concebido para colisiones de trayectoria cruzada ---ese escenario estaba fuera del alcance validado desde 2011. La filosofía moral llevaba siglos debatiendo estos dilemas; la ingeniería de software los había finalmente materializado \citep{mishra2025moral}.

Este accidente no es una anomalía. Es la expresión concreta de una tensión filosófica irresuelta: ¿bajo qué principios debe actuar una máquina cuando todas sus opciones implican daño? La industria ha respondido, mayoritariamente, con cálculos implícitos de utilidad ---priorizando la mayoría de los casos, optimizando para el usuario promedio--- sin reconocer que esta elección constituye ya una postura ética. \citet{mougan2024kantian} evidencian que las métricas de equidad dominantes en IA adoptan supuestos utilitaristas sin declararlo, violando principios de dignidad individual que la ingeniería profesional está obligada a respetar.

El Grupo 7 sostiene la siguiente tesis: \textbf{el imperativo categórico kantiano y el cálculo utilitarista de Bentham y Mill constituyen arquitecturas éticas irreconciliables para el diseño de IA en decisiones críticas}. Su confrontación revela que ningún sistema inteligente puede ser moralmente neutro: toda función objetivo, todo conjunto de reglas de prioridad y todo umbral de decisión incorpora compromisos filosóficos que deben ser explicitados, debatidos y asumidos con responsabilidad profesional. Analizar estos tres pensadores no es un ejercicio académico abstracto; es una precondición para el diseño ético de sistemas que ya toman decisiones de vida o muerte.

El presente artículo se estructura como sigue: la Sección~\ref{sec:conceptual} sintetiza los marcos filosóficos de Kant, Bentham y Mill. La Sección~\ref{sec:aplicacion} los extrapola a la arquitectura de IA y realiza la crítica cruzada entre los tres pensadores. La Sección~\ref{sec:caso} analiza el caso Tesla Autopilot bajo cada lente filosófico. La Sección~\ref{sec:impacto} examina las implicaciones para la responsabilidad profesional del ingeniero. La Sección~\ref{sec:conclusion} sintetiza los hallazgos.
